\section*{Capítulo \ref{ch:b3:measure} --- Teoria da medida}

\begin{solution}{exo:b3:measure:1}
(a) A complementação troca os dois casos da definição. Uma união
enumerável de conjuntos enumeráveis é enumerável; se um dos membros é
coenumerável, a união é coenumerável: a estabilidade vale. Ela
contém os conjuntos unitários, e toda $\sigma$-álgebra que os contenha
contém todos os conjuntos enumeráveis (uniões enumeráveis) e seus
complementares: ela é $\sigma(\{\text{conjuntos unitários}\})$.

(b) Escreva $\mathcal B = \sigma(\text{abertos})$. Todo aberto
de $\R$ é união enumerável de intervalos abertos de
extremos racionais (em torno de cada ponto racional do aberto,
um intervalo de raio racional contido nele), de modo que os abertos $\in
\sigma(\text{intervalos abertos}) \subseteq \sigma(\text{rational
data})$. Conversões: $\intoo ab = \bigcup_n\intcc{a +
\frac1n}{b - \frac1n}$; $\intcc ab = \bigcap_n \intoo{a -
\frac1n}{b + \frac1n}$; $\intoc{-\infty}a = \bigcap_n
\intoo{-\infty}{a + \frac1n}$ e reciprocamente $\intoo ab =
\intoo{-\infty}b \setminus \intoc{-\infty}a$; semirretas racionais:
$\intoc{-\infty}a =
\bigcap_{q \in \Q,\, q > a}\intoc{-\infty}q$. Cada família
gera as outras por operações enumeráveis: todas as quatro geram
$\mathcal B$.

(c) Com extremos apenas finitos, a família nem sequer é uma
álgebra: o complementar de $\intoc01$ contém semirretas ilimitadas.
Permitindo extremos infinitos ($\intoc{-\infty}b$, $\intoo
a{+\infty}$), ela se torna uma álgebra (complementares e uniões
finitas de tais uniões são desse tipo), mas não uma $\sigma$-álgebra:
$\{0\} = \bigcap_n\intoc{-\frac1n}0$ não é união finita de
intervalos não degenerados.
\end{solution}

\begin{solution}{exo:b3:measure:2}
(a) $A \cup B = A \sqcup (B \setminus (A\cap B))$, de modo que
$\mu(A\cup B) = \mu(A) + \mu(B) - \mu(A\cap B)$ (a finitude
permite a subtração). Três conjuntos: aplique duas vezes a fórmula
para dois conjuntos,
\[
\mu(A\cup B\cup C) = \sum\mu(A) - \sum\mu(A\cap B) +
\mu(A\cap B\cap C)
\]
(somas sobre os conjuntos de índices evidentes).

(b) Para a \omterm{def:b3:measure:lebesgueouter}{medida de Lebesgue}, $A_n = [n, +\infty) \downarrow
\varnothing$, mas $\lambda(A_n) = \infty \not\to 0$.

(c) Um ponto está num intervalo de comprimento $\varepsilon$:
$\lambda(\{x\}) = 0$; a subaditividade enumerável mata os conjuntos
enumeráveis. Logo $\lambda(\Q \cap \intcc01) = 0$ e, por aditividade,
$\lambda(\intcc01\setminus\Q) = 1$: os irracionais carregam todo
o comprimento.
\end{solution}

\begin{solution}{exo:b3:measure:3}
As semirretas $\intoc{-\infty}t$ formam um $\pi$-sistema (a
interseção de duas é a menor) que gera $\mathcal B(\R)$
(\cref{exo:b3:measure:1}). Os conjuntos $P_k = \intoc{-\infty}k$
crescem para $\R$ com $\mu(P_k) \leq 1 < \infty$:
o~\cref{thm:b3:measure:uniqueness} se aplica, e $\mu = \nu$ em
$\mathcal B(\R)$. Assim, a função de distribuição determina a
\omterm{def:b3:measure:measure}{medida}.
\end{solution}

\begin{solution}{exo:b3:measure:4}
Para todo $N$, $\limsup A_n \subseteq \bigcup_{n \geq N}A_n$,
de modo que $\mu(\limsup A_n) \leq \sum_{n\geq N}\mu(A_n)$, a cauda de
uma série convergente: faça $N \to \infty$. Aplicação: com
$A_n = \{x \in \intcc01 : \abs{x - r_n} \leq 4^{-n}\}$
($r_n$ o $n$-ésimo racional), $\lambda(A_n) \leq 2\cdot4^{-n}$
é somável: $\lambda(\limsup A_n) = 0$, isto é, quase todo $x$
pertence a apenas finitos $A_n$. (E, no entanto, todo $x$ é limite
de racionais: o que importa é a \emph{velocidade} $4^{-n}$.)
\end{solution}

\begin{solution}{exo:b3:measure:5}
$K = \bigcap K_n$ é uma interseção de uniões finitas de intervalos
fechados: \omterm{def:b3:topology:compact}{compacto}. Na etapa $n$ restam $2^n$ intervalos
de comprimento comum $\ell_n \leq 2^{-n}$ (cada etapa divide ao meio e
encolhe); um intervalo $I \subseteq K$ estaria dentro de um único
intervalo da etapa $n$ para todo $n$, forçando $\ell(I) = 0$: \omterm{def:b3:topology:interior}{interior}
vazio. A \omterm{def:b3:measure:measure}{medida} retirada é $\sum_{n\geq1}2^{n-1}\cdot
4^{-n} = \frac12\sum_{n\geq1}2^{-n} = \frac12$, sendo todas as retiradas
intervalos abertos disjuntos: $\lambda(K) = \frac12$.

Variante: retirar intervalos centrais de comprimento
$\varepsilon
4^{-n}$ deixa um \omterm{def:b3:topology:compact}{compacto} de \omterm{def:b3:topology:interior}{interior} vazio $K^{(\varepsilon)}$ de
\omterm{def:b3:measure:measure}{medida} $1 - \frac\varepsilon2$. Então $\bigcup_m
K^{(1/m)}$ é \omterm{rem:b3:complete:meagre}{magro} (união enumerável de conjuntos nunca densos)
de \omterm{def:b3:measure:measure}{medida} $\geq \sup_m(1 - \frac1{2m}) = 1$: um \omterm{rem:b3:complete:meagre}{conjunto magro} de
\omterm{def:b3:measure:measure}{medida} total --- e seu complementar em $\intoo01$ é um
$G_\delta$ denso de \omterm{def:b3:measure:measure}{medida} $0$ (topologicamente gordo, metricamente
nulo). O complementar de $K^{(\varepsilon)}$ em $\intoo01$ é
aberto, denso, de \omterm{def:b3:measure:measure}{medida} $\frac\varepsilon2 < \varepsilon$.
\end{solution}

\begin{solution}{exo:b3:measure:6}
Cortando $\intoc01$ em $2^n$ transladados de $\intoc0{2^{-n}}$:
$c = 2^n\,\mu(\intoc0{2^{-n}})$, de modo que $\mu(\intoc0{2^{-n}}) =
c\,2^{-n} = c\,\lambda(\intoc0{2^{-n}})$. Por invariância
por translação e aditividade, $\mu = c\lambda$ em todo intervalo
$\intoc ab$ com $b - a$ um racional diádico e qualquer $a$; um
$\intoc ab$ geral é uma união crescente de tais intervalos
($b_k \uparrow b$ passos diádicos a partir de $a$), e a \omterm{def:b3:topology:continuity}{continuidade} por
baixo estende a igualdade. Os intervalos $\intoc ab$ formam um
$\pi$-sistema que gera $\mathcal B(\R)$, com $\intoc{-k}k
\uparrow \R$ de \omterm{def:b3:measure:measure}{medida} finita ($\mu(\intoc{-k}k) = 2kc$):
o~\cref{thm:b3:measure:uniqueness} dá $\mu = c\lambda$ em
$\mathcal B(\R)$.
\end{solution}

\begin{solution}{exo:b3:measure:7}
Por regularidade (\cref{thm:b3:measure:regularity}), escolha $K
\subseteq A \subseteq U$ com $K$ \omterm{def:b3:topology:compact}{compacto}, $U$ aberto,
$\lambda(U\setminus K) < \varepsilon$ (ambas as aproximações
a menos de $\varepsilon/2$, e $\lambda(U \setminus K) =
\lambda(U\setminus A) + \lambda(A \setminus K)$). Escreva $U$ como
união enumerável disjunta de intervalos abertos $(I_n)$ (as
componentes do aberto); o \omterm{def:b3:topology:compact}{compacto} $K$ é coberto por
finitos deles, $K \subseteq B = I_1\cup\dots\cup I_N \subseteq
U$. Então $A \setminus B \subseteq A\setminus K \subseteq
U\setminus K$ e $B \setminus A \subseteq U \setminus A
\subseteq U\setminus K$:
$\lambda(A\,\triangle\,B) \leq 2\lambda(U\setminus K)$ ---
comece com $\varepsilon/2$ para cair abaixo de $\varepsilon$.
\end{solution}

\begin{solution}{exo:b3:measure:8}
Substituindo $A$ por $A \cap [-M, M]$ de \omterm{def:b3:measure:measure}{medida} positiva (algum $M$
serve, por \omterm{def:b3:topology:continuity}{continuidade} por baixo), suponha $0 < \lambda(A) <
\infty$. Tome $U \supseteq A$ aberto com $\lambda(U) <
\frac43\lambda(A)$ e decomponha $U = \bigsqcup_nI_n$ em
intervalos abertos disjuntos: $\lambda(A) = \sum_n\lambda(A\cap
I_n)$. Se todo $n$ tivesse $\lambda(A\cap I_n) \leq
\frac34\ell(I_n)$, somando obteríamos $\lambda(A) \leq
\frac34\lambda(U) < \lambda(A)$: algum intervalo $I$ satisfaz
$\lambda(A\cap I) > \frac34\ell(I)$. Ponha $B = A \cap I$ e seja
$\abs t < \frac12\ell(I)$: tanto $B$ quanto $B + t$ estão no
intervalo $I \cup (I + t)$, de comprimento $< \frac32\ell(I)$. Se
fossem disjuntos, $\lambda(B) + \lambda(B + t) = 2\lambda(B)
> \frac32\ell(I)$ superaria a \omterm{def:b3:measure:measure}{medida} do intervalo que os contém
--- impossível. Logo $B \cap (B + t) \neq \varnothing$:
algum $x \in B$ se escreve $x = y + t$ com $y \in B$, e $t = x -
y \in A - A$.
Logo $\intoo{-\frac{\ell(I)}2}{\frac{\ell(I)}2} \subseteq A -
A$.
\end{solution}

\begin{solution}{exo:b3:measure:9}
Os transladados de Vitali $(V + q)_{q\in\Q}$ formam uma partição de $\R$ (todo
real é equivalente a exatamente um representante). Suponha que todos
os conjuntos $B_q = A \cap (V + q)$ fossem mensuráveis. Dois
elementos quaisquer de $V + q$ diferem por um irracional ou por zero (dois
representantes distintos não são equivalentes), de modo que $B_q - B_q$
encontra $\Q$ apenas em $\{0\}$: ele não contém intervalo algum, e
Steinhaus (\cref{exo:b3:measure:8}) força $\lambda(B_q) = 0$.
Então $\lambda(A) \leq \sum_{q}\lambda(B_q) = 0$, contradizendo
$\lambda(A) > 0$. Logo algum $B_q \subseteq A$ não é mensurável.
\end{solution}

\begin{solution}{exo:b3:measure:10}
\emph{Mensurável $\Rightarrow$ aproximação a $\varepsilon$}: pela
regularidade exterior (\cref{thm:b3:measure:regularity}), tome $U \supseteq A$ aberto com $\lambda(U) \leq \lambda(A) +
\varepsilon$; a mensurabilidade permite a subtração
$\lambda(U\setminus A) = \lambda(U) - \lambda(A) \leq
\varepsilon$. \emph{Versão $\varepsilon$ $\Rightarrow$
versão $G_\delta$}: tome $U_n$ com $\lambda^*(U_n\setminus
A) < \frac1n$ e $G = \bigcap U_n$: um $G_\delta$ com
$\lambda^*(G\setminus A) \leq \lambda^*(U_n\setminus A) \to 0$.
\emph{Versão $G_\delta$ $\Rightarrow$ mensurável}: $G\setminus
A$ é $\lambda^*$-nulo, logo mensurável por \omterm{def:b3:complete:complete}{completude}
(\cref{thm:b3:measure:caratheodory}), e $A = G \setminus
(G\setminus A)$ é mensurável ($G$ é boreliano). Assim, os conjuntos de Lebesgue
são exatamente ``borelianos módulo nulos''.
\end{solution}

\begin{solution}{exo:b3:measure:11}
(a) $\liminf A_n = \bigcup_N\bigcap_{n\geq N}A_n$ é uma
união crescente dos conjuntos $B_N = \bigcap_{n\geq N}A_n$,
de modo que $\mu(\liminf A_n) = \lim_N\mu(B_N)$ (\omterm{def:b3:topology:continuity}{continuidade} por
baixo); e $\mu(B_N) \leq \inf_{n \geq N}\mu(A_n)$, cujo
limite é $\liminf\mu(A_n)$. Para o $\limsup$: aplique o
mesmo aos complementares dentro do ambiente de \omterm{def:b3:measure:measure}{medida} finita
$U = \bigcup A_n$ --- a \omterm{def:b3:topology:continuity}{continuidade} por cima na sequência
decrescente $C_N = \bigcup_{n\geq N}A_n$ exige $\mu(C_1)
\leq \mu(U) < \infty$ e dá $\mu(\limsup A_n) =
\lim\mu(C_N) \geq \limsup\mu(A_n)$.

(b) O intervalo móvel $A_n = \intcc n{n+1}$: todo ponto
pertence a no máximo dois dos $A_n$ e a nenhum a partir de certa ordem,
de modo que $\limsup A_n = \varnothing$; e, no entanto, $\mu(A_n) = 1$. Assim,
$\limsup\mu(A_n) = 1 > 0 = \mu(\limsup A_n)$: sem um envelope
de \omterm{def:b3:measure:measure}{medida} finita, a segunda desigualdade de (a) falha
--- a massa escapa para o infinito, onde nenhum conjunto fixo pode
capturá-la.

(c) Se $\sum\mu(A_n) < \infty$: $\mu(C_N) \leq
\sum_{n\geq N}\mu(A_n) \to 0$ e $\limsup A_n =
\bigcap C_N$ tem \omterm{def:b3:measure:measure}{medida} $\leq \inf\mu(C_N) = 0$. Casos
monótonos: crescente é a \omterm{def:b3:topology:continuity}{continuidade} por baixo; decrescente com
$\mu(A_1) < \infty$ é a \omterm{def:b3:topology:continuity}{continuidade} por cima ---
ambas demonstradas nas propriedades básicas do~\cref{ch:b3:measure}; o
contraexemplo $A_n = \intco n\infty$ (decrescendo para
$\varnothing$ com $\mu \equiv \infty$) mostra que a finitude
é, de novo, essencial.
\end{solution}

\begin{solution}{exo:b3:measure:12}
(a) Os conjuntos $E_{k,N}$ crescem com $N$ (menos
restrições), e todo $x$ acaba satisfazendo $\abs{f_n(x)
- f(x)} \leq \frac1k$ para todo $n \geq N(x)$ (convergência
pontual): $\bigcup_NE_{k,N} = X$. \omterm{def:b3:topology:continuity}{Continuidade} por baixo:
$\mu(E_{k,N}) \to \mu(X) < \infty$, de modo que $\mu(X\setminus
E_{k,N}) \to 0$; escolha $N_k$ de acordo.

(b) Seja $A = \bigcap_kE_{k,N_k}$: $\mu(X\setminus A) \leq
\sum_k\varepsilon2^{-k} = \varepsilon$. Em $A$: para todo
$k$, todos os $n \geq N_k$ satisfazem $\sup_A\abs{f_n - f} \leq
\frac1k$ --- exatamente a convergência uniforme em $A$.

(c) Para o pico móvel, a convergência uniforme em $A$ exige
que $A$ acabe evitando todo $\intcc n{n+1}$ --- mais
precisamente, $\sup_A\abs{f_n} < \frac12$ força $A \cap \intcc
n{n+1}$ a ser vazio para $n$ grande, de modo que $X \setminus A$
contém uma cauda $\bigcup_{n\geq n_0}\intcc n{n+1}$, de
\omterm{def:b3:measure:measure}{medida} infinita. Em (a), a finitude converteu
``$E_{k,N}\nearrow X$'' em ``as \omterm{def:b3:measure:measure}{medidas} dos complementares
tendem a $0$'': a \omterm{def:b3:topology:continuity}{continuidade} por cima precisa de um início finito,
e, em \omterm{def:b3:measure:measure}{espaços de medida} infinita, a fuga para o infinito é
precisamente o que ela não consegue ver.
\end{solution}

\begin{solution}{pb:b3:measure:1}
\textbf{1.} Indução. \omterm{def:b3:topology:continuity}{Continuidade}: as três fórmulas coincidem nas
junções ($\frac12c_n(1) = \frac12$ e $\frac12 +
\frac12c_n(0) = \frac12$); cada peça é \omterm{def:b3:topology:continuity}{contínua}.
A monotonicidade e os valores nos extremos são herdados. Para a
estimativa de contração: em $[0,\frac13]$, $\abs{c_{n+1} - c_n}(x)
= \frac12\abs{c_n - c_{n-1}}(3x) \leq \frac12\norm{c_n -
c_{n-1}}_\infty$; no terço médio a diferença é $0$; no
terço direito, o mesmo que no esquerdo.

\textbf{2.} $\norm{c_{n+1} - c_n}_\infty \leq
2^{-n}\norm{c_1 - c_0}_\infty$: a série dos incrementos
converge uniformemente, de modo que $c_n \to c$ uniformemente; $c$ é
\omterm{def:b3:topology:continuity}{contínua}, não decrescente, $c(0) = 0$, $c(1) = 1$ (tudo
preservado por limites uniformes) e, passando ao limite na
recursão que a define, vê-se que $c$ satisfaz ela própria as três
identidades autossemelhantes.

\textbf{3.} Pela identidade do meio, $c \equiv \frac12$ em
$\intcc{\frac13}{\frac23}$, a primeira lacuna. Toda lacuna de $C$ é
imagem da primeira lacuna por uma composição das duas
contrações afins $x \mapsto \frac x3$, $x\mapsto\frac{x +
2}3$; as identidades transportam a constância correspondentemente (com
valores nos racionais diádicos). Para a fórmula dos dígitos, tome $x =
\sum_n 2b_n3^{-n} \in C$: se $b_1 = 0$, então $x \in
[0,\frac13]$ e $c(x) = \frac12c(3x)$, com $3x$ tendo dígitos
$(b_2, b_3, \dots)$; se $b_1 = 1$, então $x \in [\frac23, 1]$
e $c(x) = \frac12 + \frac12c(3x - 2)$, mesmo deslocamento. Por
indução, os $N$ primeiros dígitos binários de $c(x)$ são $b_1,
\dots, b_N$ para todo $N$: $c(x) = \sum_nb_n2^{-n}$.

\textbf{4.} Fora de $C$, $c$ é localmente constante: derivável,
com derivada $0$. Como $\lambda(C) = 0$
(\cref{ex:b3:measure:cantor}), $c' = 0$ quase em toda parte. E, no entanto,
$c(1) - c(0) = 1$: o teorema fundamental na forma $\mathcal
C^1$ exige que $c$ seja derivável \emph{em toda parte},
com derivada \omterm{def:b3:topology:continuity}{contínua} (ou ao menos integrável, mais a \omterm{def:b3:topology:continuity}{continuidade}
absoluta --- veja o~\cref{ch:b3:lebesgue}); $c$ não é
derivável nos pontos de $C$ e, mais fundamentalmente, $c$
falha na \emph{\omterm{def:b3:topology:continuity}{continuidade} absoluta}: ela sobe num conjunto nulo.

\textbf{5.} Dado $y = \sum_n\beta_n2^{-n} \in \intcc01$
($\beta_n \in \{0,1\}$), o ponto $x = \sum_n2\beta_n3^{-n}
\in C$ tem $c(x) = y$ pela questão 3: $c(C) = \intcc01$, um conjunto
de \omterm{def:b3:measure:measure}{medida} $1$ --- o conjunto nulo $C$ carrega, através de $c$, o
intervalo inteiro.

\textbf{6.} $h$ é \omterm{def:b3:topology:continuity}{contínua} e estritamente crescente ($x$ é,
$c$ é não decrescente); $h(0) = 0$, $h(1) = 1$, de modo que, pelo
teorema do valor intermediário, $h$ é uma bijeção \omterm{def:b3:topology:continuity}{contínua} de
$\intcc01$; uma bijeção \omterm{def:b3:topology:continuity}{contínua} de um \omterm{def:b3:topology:compact}{espaço compacto} em um \omterm{def:b3:topology:hausdorff}{de
Hausdorff} é um \omterm{def:b3:topology:continuity}{homeomorfismo}
(\cref{cor:b3:topology:compacthomeo}).

\textbf{7.} Numa lacuna $(u, v)$ (comprimento $\ell$), $c$ é constante,
de modo que $h$ é afim de inclinação $\frac12$: $h((u,v))$ é um
intervalo de comprimento $\frac\ell2$. As lacunas são disjuntas e $h$
é injetora: as imagens são disjuntas, de \omterm{def:b3:measure:measure}{medida} total
$\frac12\sum\ell = \frac12(1 - \lambda(C)) = \frac12$.

\textbf{8.} $h(\intcc01) = \intcc01$ e $h(C)$ é \omterm{def:b3:topology:compact}{compacto}
(imagem \omterm{def:b3:topology:continuity}{contínua}), logo mensurável, com
\[
\lambda\bigl(h(C)\bigr) = 1 -
\lambda\bigl(h(\intcc01\setminus C)\bigr) = 1 - \tfrac12 =
\tfrac12 .
\]
Um \omterm{def:b3:topology:continuity}{homeomorfismo} pode inflar um conjunto nulo até a \omterm{def:b3:measure:measure}{medida} $\frac12$:
``tamanho \omterm{def:b3:topology:topology}{topológico}'' (categoria, dimensão) e ``\omterm{def:b3:measure:measure}{medida}'' são
transportados por \omterm{def:b3:topology:continuity}{homeomorfismos} de maneiras muito diferentes --- só o
primeiro é um invariante \omterm{def:b3:topology:topology}{topológico}.

\textbf{9.} $Z = h^{-1}(W) \subseteq h^{-1}(h(C)) = C$ tem
$\lambda^*(Z) \leq \lambda(C) = 0$: um conjunto nulo, logo
Lebesgue-mensurável por \omterm{def:b3:complete:complete}{completude}
(\cref{thm:b3:measure:caratheodory}).

\textbf{10.} Seja $\varphi$ \omterm{def:b3:topology:continuity}{contínua} e $\mathcal D = \{B
: \varphi^{-1}(B) \in \mathcal B\}$. As pré-imagens comutam com
complementares e uniões enumeráveis, de modo que $\mathcal D$ é uma
$\sigma$-álgebra; ela contém os abertos (\omterm{def:b3:topology:continuity}{continuidade}):
$\mathcal D \supseteq \mathcal B$ --- pré-imagens de borelianos
por \omterm{def:b3:topology:continuity}{aplicações contínuas} são borelianas.

\textbf{11.} Se $Z$ fosse boreliano, aplique a questão 10 à
\omterm{def:b3:topology:continuity}{contínua} $\varphi = h^{-1}$: $\varphi^{-1}(Z) = h(Z) = W$
seria boreliano, logo Lebesgue-mensurável --- contradizendo a
escolha de $W$. Logo $Z \in \mathcal L \setminus \mathcal B$:
a $\sigma$-álgebra de Lebesgue contém estritamente a de Borel.

\textbf{12.} $g = \mathbf 1_Z$ é Lebesgue-mensurável ($Z \in
\mathcal L$) e $\varphi = h^{-1}$ é \omterm{def:b3:topology:continuity}{contínua}, mas
$(g\circ\varphi)^{-1}(\{1\}) = \varphi^{-1}(Z) = W$ não é
mensurável: $g \circ \varphi$ não é Lebesgue-mensurável. O
cuidado necessário: ``função Lebesgue-mensurável'' significa que as pré-imagens
de conjuntos \emph{borelianos} caem em $\mathcal L$; compor exige que as
pré-imagens de conjuntos \emph{de Lebesgue} sejam de Lebesgue, o que
a \omterm{def:b3:topology:continuity}{continuidade} não garante (aqui $\varphi^{-1}(Z) \notin
\mathcal L$, ainda que $\varphi$ seja um \omterm{def:b3:topology:continuity}{homeomorfismo}).

\textbf{13.} As cadeias, com testemunhas da estrita inclusão:
\[
\{\text{enumerável}\} \subsetneq \{\text{Borel-nulo}\}
\subsetneq \{\text{Lebesgue-nulo}\} \subsetneq \mathcal L
\subsetneq \mathcal P(\R),
\qquad
\{\text{Borel-nulo}\} \subsetneq \mathcal B \subsetneq
\mathcal L .
\]
Testemunhas: $C$ é boreliano, nulo, não enumerável (primeira lacuna); $Z$ é
Lebesgue-nulo, mas não boreliano (segunda e, dentro de $\mathcal B
\subsetneq \mathcal L$, sexta); um conjunto de Cantor gordo é boreliano,
nunca denso, de \omterm{def:b3:measure:measure}{medida} positiva (separando os conjuntos nulos dos
borelianos); o $V$ de Vitali não está em $\mathcal L$ (última lacuna).
\omterm{def:b3:measure:measure}{Medida}, \omterm{def:b3:topology:topology}{topologia} e cardinalidade fatiam $\mathcal P(\R)$ segundo
linhas genuinamente diferentes.

\textbf{14.} \omterm{def:b3:measure:outer}{Medida exterior}: $\mu_F^*(\varnothing) = 0$ (cubra
por um intervalo que se anula), a monotonicidade é clara, e a subaditividade
enumerável decorre de emendar coberturas $\varepsilon2^{-k}$-ótimas,
exatamente como para $\lambda^*$. A cobertura por um único intervalo
dá $\mu_F^*(\intoc ab) \leq F(b) - F(a)$. Reciprocamente, seja
$\intoc ab \subseteq \bigcup_k\intoc{a_k}{b_k}$. Pela \omterm{def:b3:topology:continuity}{continuidade}
de $F$, tome $b_k' > b_k$ com $F(b_k') \leq F(b_k) +
\varepsilon2^{-k}$, e $a' \in \intoo ab$ com $F(a') \leq
F(a) + \varepsilon$. O \omterm{def:b3:topology:compact}{compacto} $\intcc{a'}b$ é coberto pelos
abertos $\intoo{a_k}{b_k'}$: finitos bastam, e o
argumento de encadeamento do~\cref{thm:b3:measure:lebesgue} (caminhe de
$a'$ até $b$ por intervalos que se sobrepõem, telescopando
os incrementos de $F$, com a monotonicidade absorvendo as sobreposições) fornece $F(b)
- F(a') \leq \sum_k(F(b_k') - F(a_k)) \leq \sum_k(F(b_k) -
F(a_k)) + \varepsilon$. Faça $\varepsilon \to 0$: igualdade.
(A \omterm{def:b3:topology:continuity}{continuidade} de $F$ é o que permitiu abrir os intervalos a um
custo em $F$ arbitrariamente pequeno.)

\textbf{15.} Basta demonstrar que cada semirreta $H_t =
\intoc{-\infty}t$ é Carathéodory-mensurável, já que os conjuntos
mensuráveis formam uma $\sigma$-álgebra
(\cref{thm:b3:measure:caratheodory}) e as semirretas geram
$\mathcal B$. Dados $A$ e uma cobertura $\varepsilon$-ótima
$(\intoc{a_k}{b_k})$ de $A$: cada intervalo se reparte como
$\intoc{a_k}{t\wedge b_k} \cup \intoc{t \vee a_k}{b_k}$ (uma das peças
possivelmente vazia), com custos em $F$ que somam exatamente
$F(b_k) - F(a_k)$; as primeiras peças cobrem $A \cap H_t$, as
segundas $A \setminus H_t$. Logo $\mu_F^*(A\cap H_t) +
\mu_F^*(A\setminus H_t) \leq \mu_F^*(A) + \varepsilon$, e
a desigualdade recíproca é a subaditividade. Restringindo a
\omterm{def:b3:measure:measure}{medida} resultante a $\mathcal B$: a \omterm{def:b3:measure:measure}{medida} de Lebesgue--Stieltjes
$\mu_F$.

\textbf{16.} $\mu(\R) = \lim_n(F(n) - F(-n)) = 1 - 0 = 1$
(\omterm{def:b3:topology:continuity}{continuidade} da \omterm{def:b3:measure:measure}{medida} ao longo de $\intoc{-n}n$). Numa lacuna
$\intoo uv$ de $C$, $c$ é constante, de modo que todo subintervalo
semiaberto é $\mu$-nulo, e a lacuna também (união
enumerável); fora de $\intcc01$, $c$ também é constante. Logo
$\mu(\R\setminus C) = 0$, $\mu(C) = 1$, ao passo que $\lambda(C) =
0$ (\cref{ex:b3:measure:cantor}): cada uma de $\mu, \lambda$ é
carregada por um conjunto que a outra declara nulo --- mutuamente
singulares.

\textbf{17.} $\mu(\{x\}) = \lim_{\delta\downarrow0}
\mu(\intoc{x-\delta}x) = \lim(c(x) - c(x-\delta)) = 0$ pela
\omterm{def:b3:topology:continuity}{continuidade} de $c$: sem átomos. Logo $\mu$ é uma \omterm{def:b3:measure:measure}{medida} de
probabilidade sem átomos carregada por um \omterm{def:b3:topology:compact}{compacto} Lebesgue-nulo ---
nem difusa com densidade como as restrições de $\lambda$,
nem atômica como as \omterm{def:b3:measure:measure}{medidas} de contagem: uma terceira espécie.

\textbf{18.} A peça $C_\varepsilon$ ocupa um intervalo
ternário $I_\varepsilon$ de comprimento $3^{-m}$, e a questão 3
da Parte I mostra que, ao longo de $I_\varepsilon$, a escada
sobe exatamente $2^{-m}$ (os $m$ primeiros dígitos binários de $c$
ficam congelados em $\varepsilon$, e os demais varrem tudo).
Logo $\mu(C_\varepsilon) = \mu(I_\varepsilon) =
c(\text{extremo direito}) - c(\text{extremo esquerdo}) = 2^{-m}$: os cilindros de dígitos de
profundidade $m$ têm todos massa $2^{-m}$, a lei de $m$ moedas
honestas.

\textbf{19.} O membro direito define a \omterm{def:b3:measure:measure}{medida} boreliana
$\nu = \frac12\,(x \mapsto \tfrac x3)_*\mu + \frac12\,(x
\mapsto \tfrac{x+2}3)_*\mu$ avaliada em $A$ --- \omterm{def:b3:measure:measure}{medida} de
probabilidade. Pontualmente, com $c$ estendida globalmente, verifica-se
caso a caso ($x \leq 0$; os três terços; $x \geq 1$) a
identidade
\[
c(x) = \tfrac12\,c(3x) + \tfrac12\,c(3x - 2),
\]
por exemplo, em $\intcc{1/3}{2/3}$: $\frac12\cdot1 + \frac12\cdot0 =
\frac12 = c(x)$. Avaliar $\nu$ em $\intoc ab$ dá, portanto,
$c(b) - c(a) = \mu(\intoc ab)$, e duas \omterm{def:b3:measure:measure}{medidas} finitas que
coincidem no $\pi$-sistema dos intervalos semiabertos coincidem
em $\mathcal B$ (\cref{thm:b3:measure:uniqueness}): $\nu =
\mu$.

\textbf{20.} Por indução sobre $n$: $c_0(1-x) = 1 - c_0(x)$
e, se $c_n$ tem a simetria, então, para $x \in
\intcc0{1/3}$: $c_{n+1}(1 - x) = \frac12 + \frac12c_n(3(1-x)
- 2) = \frac12 + \frac12c_n(1 - 3x) = \frac12 + \frac12(1 -
c_n(3x)) = 1 - c_{n+1}(x)$; o terço médio é espelhado em torno de
$\frac12$; o terço direito é o caso esquerdo refletido. No
limite, $c(1-x) = 1 - c(x)$. Imagem direta: $(s_*\mu)(\intoc ab)
= \mu(\intco{1-b}{1-a}) = c(1-a) - c(1-b)$ ($\mu$ sem átomos,
questão 17, de modo que as convenções de bordo nada custam) $= (1 -
c(a)) - (1 - c(b)) = \mu(\intoc ab)$: $s_*\mu = \mu$ por
unicidade.

\textbf{21.} Seja $X \sim \mu$ (as integrais de funções
\omterm{def:b3:topology:continuity}{contínuas} contra $\mu$ existem como limites de somas do tipo de Riemann
sobre as peças de profundidade $m$, cada uma de massa $2^{-m}$, com
erro de amostragem $\leq \operatorname{osc} \leq
\norm{f'}_\infty3^{-m}$; o~\cref{ch:b3:lebesgue}
sistematizará isso). Simetria: $1 - X \sim X$, de modo que $\E X =
\frac12$. Autossemelhança: $X$ tem a lei de $\frac Y3$ com
probabilidade $\frac12$ e a de $\frac{Y+2}3$ com probabilidade
$\frac12$, $Y \sim \mu$, de modo que
\[
\E X^2 = \frac12\,\frac{\E Y^2}9 + \frac12\,
\frac{\E Y^2 + 4\E Y + 4}{9} = \frac{2\E X^2 + 6}{18},
\]
donde $\E X^2 = \frac38$ e $\operatorname{Var}X = \frac38
- \frac14 = \frac18$. A lei uniforme tem variância
$\frac1{12} < \frac18$: a massa de Cantor se agarra aos extremos.

\textbf{22.} $C$ é fechado e $\mu(C) = 1$. Se um aberto $I$
encontra $C$ em $x$, as peças de profundidade $m$ que contêm $x$ encolhem
para $x$, de modo que alguma $I_\varepsilon \subseteq I$ e $\mu(I) \geq
2^{-m} > 0$: nenhum fechado menor pode carregar $\mu$. Pontos
fora de $C$ têm \omterm{def:b3:topology:topology}{vizinhanças} em lacunas de $\mu$-medida $0$. Logo
$\operatorname{supp}\mu = C$ exatamente.

\textbf{23.} Um fenômeno, dois dialetos. A Parte I diz que o
crescimento de $c$ é invisível à sua derivada: $c' = 0$ q.t.p.,
toda a subida concentrada no conjunto nulo $C$. A Parte V
diz que a \omterm{def:b3:measure:measure}{medida} associada $\mu_c$ põe toda a sua massa nesse
mesmo conjunto nulo: $\mu_c \perp \lambda$, de modo que nenhuma densidade $f \geq
0$ pode satisfazer $\mu_c(A) = \int_Af\,\dd\lambda$ --- uma densidade
obriga a anulação em conjuntos $\lambda$-nulos. O dicionário:
$F \leftrightarrow$ \omterm{def:b3:measure:measure}{medida} finita $\mu_F$ para toda função limitada não
decrescente (questões 14--15); $F$ é integral de sua derivada
$\leftrightarrow$ $\mu_F$ tem densidade (o caso ``absolutamente
\omterm{def:b3:topology:continuity}{contínuo}''); e, em geral, $F'$, que existe q.t.p.\ para
$F$ monótona (teorema de derivação de Lebesgue, além
deste capítulo), recupera apenas a parte com densidade. A escada
é o extremo: \omterm{def:b3:topology:continuity}{contínua}, com derivada $0$ q.t.p.\ ---
sua \omterm{def:b3:measure:measure}{medida} é \emph{puramente singular} e o teorema fundamental
do cálculo, longe de falhar por acidente, falha exatamente pela
quantidade $\mu_c(\R) = 1$ de massa singular.

\textbf{24.} Fixe $x < y$ em $\intcc01$ e escolha $m \geq 0$
com $3^{-(m+1)} < y - x \leq 3^{-m}$. Ao longo de qualquer intervalo
triádico $\intcc{k3^{-m}}{(k+1)3^{-m}}$, a escada sobe
no máximo $2^{-m}$: um tal intervalo ou é uma das $2^m$
peças de $C_m$, em que a subida é exatamente $2^{-m}$ (questão
18), ou está contido no \omterm{def:b3:topology:interior}{fecho} de uma única lacuna de alguma
etapa $\leq m$, em que $c$ é constante. Como $y - x \leq
3^{-m}$, o intervalo $\intcc xy$ encontra no máximo dois
intervalos triádicos consecutivos de profundidade $m$, de modo que
\[
c(y) - c(x) \leq 2\cdot2^{-m}
= 2\bigl(3^{-m}\bigr)^{s}
< 2\bigl(3(y - x)\bigr)^{s}
= 4\,(y - x)^{s},
\]
usando $3^{s} = 2$ e $3^{-m} < 3(y - x)$. Otimalidade: os
extremos $u < v$ de uma peça de Cantor de profundidade $m$ satisfazem $v - u =
3^{-m}$ e $c(v) - c(u) = 2^{-m} = (v - u)^{s}$; uma cota de Hölder
$\abs{c(v) - c(u)} \leq K(v - u)^{t}$ com $t > s$ forçaria
$K \geq 2^{-m}3^{mt} = (3^{t}/2)^{m} \to \infty$ (pois
$3^{t} > 3^{s} = 2$), e todo subintervalo de $\intcc01$
contém tais peças, de modo que a falha é local também.
Forma em \omterm{def:b3:measure:measure}{medida}: $\mu\bigl(\intcc{x-r}{x+r}\bigr) \leq
c(x + r) - c(x - r) \leq 4(2r)^{s} = 4\cdot2^{s}r^{s} \leq
8r^{s}$ (valores de $c$ estendidos a $\R$ como na questão 16;
$\mu$ não tem átomos, questão 17).

\textbf{25.} Escreva $S_0(x) = \frac x3$ e $S_1(x) =
\frac{x+2}3$; a hipótese diz que $\nu = \frac12(S_0)_*\nu +
\frac12(S_1)_*\nu$. Iterando $m$ vezes,
\[
\nu = 2^{-m}\sum_{w \in \{0,1\}^m}(S_w)_*\nu,
\qquad S_w = S_{w_1}\circ\dots\circ S_{w_m}.
\]
Como $\nu$ é carregada por $\intcc01$ e $S_w(\intcc01) =
I_w$, a peça de Cantor de profundidade $m$ indexada por $w$, cada
$(S_w)_*\nu$ é uma \omterm{def:b3:measure:measure}{medida} de probabilidade carregada por $I_w$; as
$2^m$ peças são intervalos fechados dois a dois disjuntos de comprimento
$3^{-m}$. Fixe $\intoc ab$ e seja $N_m$ o número de
peças $I_w \subseteq \intoc ab$. Uma peça que encontra $\intoc ab$
sem estar contida nele deve conter $a$ ou $b$, e um
ponto está em no máximo uma peça, de modo que
\[
2^{-m}N_m \leq \nu(\intoc ab) \leq 2^{-m}N_m + 2\cdot2^{-m}.
\]
A mesma dupla desigualdade vale para $\mu$ (a questão 19 dá
a iteração idêntica), com a mesma contagem $N_m$. Logo
$\abs{\nu(\intoc ab) - \mu(\intoc ab)} \leq 2^{-m+1} \to 0$:
$\nu$ e $\mu$ coincidem no $\pi$-sistema dos intervalos
semiabertos, e ambas são \omterm{def:b3:measure:measure}{medidas} de probabilidade, de modo que
o~\cref{thm:b3:measure:uniqueness} dá $\nu = \mu$. A
\omterm{def:b3:measure:measure}{medida} da escada é, assim, \emph{o} ponto fixo do esquema de
média com duas aplicações --- o enunciado, no nível das \omterm{def:b3:measure:measure}{medidas},
da autossemelhança de $C$.
\end{solution}
